% GENERATED FILE. Edit canonical lessons, then run npm run generate:books.
% canonical-source-hash: fnv1a64:f1b2e33c448dcc60
% canonical-lessons: KA-C471-gaddala, KA-C471-jamayisu, KA-C471-dittisu, KA-C471-jaggu, KA-C471-tuluku

\chapter{Crowds and Noise}
\label{ch:ka-crowds-and-noise}
\hlchaptermodality{471}

\begin{chapteropening}
\textbf{By the end of this chapter:} \emph{I can say a din, to gather, to stare, to tug and to spill over.}

Complete the last lesson of chapter 471: I can say a din, to gather, to stare, to tug and to spill over.
\end{chapteropening}

\section[\texorpdfstring{\kn{ಗದ್ದಲ} (gaddala)}{ಗದ್ದಲ (gaddala)}]{\kn{ಗದ್ದಲ} (gaddala) — a din; a commotion}

\label{lesson:KA-C471-gaddala}



\begin{quote}
Before the new one: say the Kannada for supplies, then the Kannada for a telephone.
\end{quote}

\subsection*{You'll want to know: \kn{ಗದ್ದಲ}}
\textbf{\kn{ಗದ್ದಲ}} — \emph{gaddala} — \textquotedblleft{}a din; a commotion\textquotedblright{}.

\begin{quote}
\kn{ಮಾರುಕಟ್ಟೆಯಲ್ಲಿ} \kn{ತುಂಬಾ} \kn{ಗದ್ದಲ} \kn{ಇತ್ತು}.
\end{quote}

\begin{quote}
\emph{mārukaṭṭeyalli tumbā gaddala ittu.}
\end{quote}

\begin{quote}
There was a lot of noise in the market.
\end{quote}

Many voices at once; \kn{ಗಲಾಟೆ} is a near twin.

\begin{cousinweb}
\textbf{\kn{ಗದ್ದಲ}} (\emph{gaddala}) is usually traced to the native family of \emph{gadaṟu}, \emph{gaddisu}, \textquotedblleft{}to make a loud sound, to roar, to shout\textquotedblright{} (DEDR 1189).
\end{cousinweb}

\begin{grammarlens}[title={What you've built}]
One more word for getting around.
\end{grammarlens}

\subsection*{Your turn}
\begin{itemize}
\raggedright
  \item \emph{Say it:} \emph{gaddala}
  \item \emph{Say it:} \emph{gaddala}, once more
  \item \emph{Say it:} \emph{mārukaṭṭeyalli tumbā gaddala ittu.} — There was a lot of noise in the market.
  \item \emph{Recall:} say the Kannada for supplies, then the Kannada for a telephone, then say \emph{gaddala} again
\end{itemize}

\begin{tcolorbox}[breakable,colback=teal!4,colframe=teal!35!black,title={Before you move on}]
What does \kn{ಗದ್ದಲ} mean? (\textquotedblleft{}A din; a commotion\textquotedblright{}.) Say it once more.
\end{tcolorbox}
\section[\texorpdfstring{\kn{ಜಮಾಯಿಸು} (jamāyisu)}{ಜಮಾಯಿಸು (jamāyisu)}]{\kn{ಜಮಾಯಿಸು} (jamāyisu) — to gather; to assemble}

\label{lesson:KA-C471-jamayisu}



\begin{quote}
Before the new one: say the Kannada for a telephone, then the Kannada for a din.
\end{quote}

\subsection*{You'll want to know: \kn{ಜಮಾಯಿಸು}}
\textbf{\kn{ಜಮಾಯಿಸು}} — \emph{jamāyisu} — \textquotedblleft{}to gather; to assemble\textquotedblright{}.

\begin{quote}
\kn{ಅಪಘಾತದ} \kn{ನಂತರ} \kn{ರಸ್ತೆಯಲ್ಲಿ} \kn{ಜನ} \kn{ಜಮಾಯಿಸಿದರು}.
\end{quote}

\begin{quote}
\emph{apaghātada nantara rasteyalli jana jamāyisidaru.}
\end{quote}

\begin{quote}
After the accident people gathered on the road.
\end{quote}

Past (I …): \kn{ಜಮಾಯಿಸಿದೆ}, \emph{jamāyiside}.

Said of a crowd that forms.

\begin{cousinweb}
\textbf{\kn{ಜಮಾಯಿಸು}} (\emph{jamāyisu}) is usually traced to an Arabic word for gathering, through Hindustani \emph{jamā}, \textquotedblleft{}gathered, collected\textquotedblright{}, with the Kannada verb ending -\emph{isu}. \kn{ಜಮಾ} \kn{ಮಾಡು}, to deposit, comes from the same word.
\end{cousinweb}

\begin{grammarlens}[title={What you've built}]
One more word for getting around.
\end{grammarlens}

\subsection*{Your turn}
\begin{itemize}
\raggedright
  \item \emph{Say it:} \emph{jamāyisu}
  \item \emph{Say it:} \emph{jamāyisu}, once more
  \item \emph{Say it:} \emph{apaghātada nantara rasteyalli jana jamāyisidaru.} — After the accident people gathered on the road.
  \item \emph{Recall:} say the Kannada for a telephone, then the Kannada for a din, then say \emph{jamāyisu} again
\end{itemize}

\begin{tcolorbox}[breakable,colback=teal!4,colframe=teal!35!black,title={Before you move on}]
What does \kn{ಜಮಾಯಿಸು} mean? (\textquotedblleft{}To gather; to assemble\textquotedblright{}.) Say it once more.
\end{tcolorbox}
\section[\texorpdfstring{\kn{ದಿಟ್ಟಿಸು} (diṭṭisu)}{ದಿಟ್ಟಿಸು (diṭṭisu)}]{\kn{ದಿಟ್ಟಿಸು} (diṭṭisu) — to stare; to gaze}

\label{lesson:KA-C471-dittisu}



\begin{quote}
Before the new one: say the Kannada for a din, then the Kannada for to gather.
\end{quote}

\subsection*{You'll want to know: \kn{ದಿಟ್ಟಿಸು}}
\textbf{\kn{ದಿಟ್ಟಿಸು}} — \emph{diṭṭisu} — \textquotedblleft{}to stare; to gaze\textquotedblright{}.

\begin{quote}
\kn{ಆ} \kn{ಅಪರಿಚಿತ} \kn{ನನ್ನನ್ನೇ} \kn{ದಿಟ್ಟಿಸುತ್ತಿದ್ದನು}.
\end{quote}

\begin{quote}
\emph{ā aparicita nannannē diṭṭisuttiddanu.}
\end{quote}

\begin{quote}
That stranger kept staring at me.
\end{quote}

Past (I …): \kn{ದಿಟ್ಟಿಸಿದೆ}, \emph{diṭṭiside}.

The one stared at takes -\kn{ಅನ್ನು}.

\begin{cousinweb}
\textbf{\kn{ದಿಟ್ಟಿಸು}} (\emph{diṭṭisu}) is built on \emph{diṭṭi}, \textquotedblleft{}a gaze\textquotedblright{}, usually traced to Sanskrit \emph{dṛṣṭi}, \textquotedblleft{}seeing, sight\textquotedblright{} (CDIAL 6520), through Prakrit, with the verb ending -\emph{isu}.
\end{cousinweb}

\begin{grammarlens}[title={What you've built}]
One more word for getting around.
\end{grammarlens}

\subsection*{Your turn}
\begin{itemize}
\raggedright
  \item \emph{Say it:} \emph{diṭṭisu}
  \item \emph{Say it:} \emph{diṭṭisu}, once more
  \item \emph{Say it:} \emph{ā aparicita nannannē diṭṭisuttiddanu.} — That stranger kept staring at me.
  \item \emph{Recall:} say the Kannada for a din, then the Kannada for to gather, then say \emph{diṭṭisu} again
\end{itemize}

\begin{tcolorbox}[breakable,colback=teal!4,colframe=teal!35!black,title={Before you move on}]
What does \kn{ದಿಟ್ಟಿಸು} mean? (\textquotedblleft{}To stare; to gaze\textquotedblright{}.) Say it once more.
\end{tcolorbox}
\section[\texorpdfstring{\kn{ಜಗ್ಗು} (jaggu)}{ಜಗ್ಗು (jaggu)}]{\kn{ಜಗ್ಗು} (jaggu) — to tug; to pull hard}

\label{lesson:KA-C471-jaggu}



\begin{quote}
Before the new one: say the Kannada for to gather, then the Kannada for to stare.
\end{quote}

\subsection*{You'll want to know: \kn{ಜಗ್ಗು}}
\textbf{\kn{ಜಗ್ಗು}} — \emph{jaggu} — \textquotedblleft{}to tug; to pull hard\textquotedblright{}.

\begin{quote}
\kn{ಗುಂಪಿನಲ್ಲಿ} \kn{ಮಗು} \kn{ನನ್ನ} \kn{ಕೈ} \kn{ಜಗ್ಗಿತು}.
\end{quote}

\begin{quote}
\emph{gumpinalli magu nanna kai jaggitu.}
\end{quote}

\begin{quote}
In the crowd the child tugged at my hand.
\end{quote}

Past (I …): \kn{ಜಗ್ಗಿದೆ}, \emph{jaggide}.

Also to sag under weight: \kn{ಹಗ್ಗ} \kn{ಜಗ್ಗಿತು}.

\begin{cousinweb}
\textbf{\kn{ಜಗ್ಗು}} (\emph{jaggu}) is native Dravidian: DEDR 2360 lists Kannada \emph{jaggu}, \textquotedblleft{}to bend down, to sink under a weight\textquotedblright{}, in the family of \emph{sari}, \textquotedblleft{}to slide\textquotedblright{}. How the sense of tugging grew from it is not settled.
\end{cousinweb}

\begin{grammarlens}[title={What you've built}]
One more word for getting around.
\end{grammarlens}

\subsection*{Your turn}
\begin{itemize}
\raggedright
  \item \emph{Say it:} \emph{jaggu}
  \item \emph{Say it:} \emph{jaggu}, once more
  \item \emph{Say it:} \emph{gumpinalli magu nanna kai jaggitu.} — In the crowd the child tugged at my hand.
  \item \emph{Recall:} say the Kannada for to gather, then the Kannada for to stare, then say \emph{jaggu} again
\end{itemize}

\begin{tcolorbox}[breakable,colback=teal!4,colframe=teal!35!black,title={Before you move on}]
What does \kn{ಜಗ್ಗು} mean? (\textquotedblleft{}To tug; to pull hard\textquotedblright{}.) Say it once more.
\end{tcolorbox}
\section[\texorpdfstring{\kn{ತುಳುಕು} (tuḷuku)}{ತುಳುಕು (tuḷuku)}]{\kn{ತುಳುಕು} (tuḷuku) — to spill over; to brim over}

\label{lesson:KA-C471-tuluku}



\begin{quote}
Before the new one: say the Kannada for to stare, then the Kannada for to tug.
\end{quote}

\subsection*{You'll want to know: \kn{ತುಳುಕು}}
\textbf{\kn{ತುಳುಕು}} — \emph{tuḷuku} — \textquotedblleft{}to spill over; to brim over\textquotedblright{}.

\begin{quote}
\kn{ಹಬ್ಬದ} \kn{ದಿನ} \kn{ಬಸ್ಸುಗಳು} \kn{ಜನರಿಂದ} \kn{ತುಳುಕುತ್ತಿದ್ದವು}.
\end{quote}

\begin{quote}
\emph{habbada dina bassugaḷu janarinda tuḷukuttiddavu.}
\end{quote}

\begin{quote}
On the festival day the buses were overflowing with people.
\end{quote}

Past (it …): \kn{ತುಳುಕಿತು}, \emph{tuḷukitu}.

Of a full pot, a crowded bus or a full heart.

\begin{cousinweb}
\textbf{\kn{ತುಳುಕು}} (\emph{tuḷuku}) is native Dravidian: DEDR 3361 lists Kannada \emph{tuḷaku}, \textquotedblleft{}to be scattered in drops, to spill, to run over\textquotedblright{}, and \emph{tuḷuku}.
\end{cousinweb}

\begin{grammarlens}[title={What you've built}]
One more word for getting around.
\end{grammarlens}

\subsection*{Your turn}
\begin{itemize}
\raggedright
  \item \emph{Say it:} \emph{tuḷuku}
  \item \emph{Say it:} \emph{tuḷuku}, once more
  \item \emph{Say it:} \emph{habbada dina bassugaḷu janarinda tuḷukuttiddavu.} — On the festival day the buses were overflowing with people.
  \item \emph{Recall:} say the Kannada for to stare, then the Kannada for to tug, then say \emph{tuḷuku} again
\end{itemize}

\begin{tcolorbox}[breakable,colback=teal!4,colframe=teal!35!black,title={Before you move on}]
What does \kn{ತುಳುಕು} mean? (\textquotedblleft{}To spill over; to brim over\textquotedblright{}.) Say it once more.
\end{tcolorbox}
